(1) 单位宽度流率 \(q=u_1h_1=u_2h_2\)。对控制体积分压强的静水力和动量通量得 \(q^2/h_1+gh_1^2/2=q^2/h_2+gh_2^2/2\)。无须在薄水跃内假定流体无耗散，因为这里保留的是总水平动量。

(2) 设 \(r=h_2/h_1\)，把 \(q^2=Fr_1^2g h_1^3\) 代入、消去平凡根 \(r=1\)，得到 \(r^2+r=2Fr_1^2\)，即 \(r=[\sqrt{1+8Fr_1^2}-1]/2>1\)。下游 \(Fr_2^2=Fr_1^2/r^3=(r+1)/(2r^2)<1\)（当 \(r>1\) 时）。

(3) 用 \(u_2=u_1/r\) 有 \(E_1-E_2=h_1(1-r)+\frac{Fr_1^2h_1}{2}(1-r^{-2})\)。代入 \(Fr_1^2=r(r+1)/2\) 后化简为 \(\boxed{E_1-E_2=(h_2-h_1)^3/(4h_1h_2)}\)，仅当无水跃 \(h_2=h_1\) 时为零。真实能量损失转为湍流和热。

(4) 若 \(Fr_1=1+\epsilon\)，由 \(r^2+r=2(1+2\epsilon+\cdots)\) 得 \(r=1+4\epsilon/3+O(\epsilon^2)\)，所以比能损失为 \(O(\epsilon^3)\)。强水跃 \(Fr_1\gg1\) 时 \(r\sim\sqrt2 Fr_1\)，\(E_1-E_2\sim h_2^2/(4h_1)\)。水跃产生不可逆熵增，而零耗散的线性浅水波无该有限跳变。
